\documentclass[10pt,a4paper]{amsart}
\usepackage[margin=19mm]{geometry}
\usepackage{amsmath,amssymb,mathtools}
\usepackage[scr=rsfs,cal=euler]{mathalfa}
\usepackage{xfrac}
\usepackage[unicode,hidelinks,bookmarks=false]{hyperref}
\newcommand{\ie}{i.e.\ }
\newcommand{\cf}{cf.\ }
\newcommand{\resp}{respectively\ }
\newcommand{\Subsection}{\subsection}
\providecommand{\nobreakdash}{\nobreak\hbox{-}}
\setlength{\emergencystretch}{2em}
\multlinegap0pt
\usepackage{amssymb}
\usepackage{enumerate}
\usepackage{bm}
\newtheorem{prop}{Proposition}
\title{The Zero Set of an Electric Field from a Finite Number of Point Charges}
\author{Tam\'as Erd\'elyi, Joseph Rosenblatt, Rebecca Rosenblatt}
\date{Source: arXiv:2106.04706v2 (CC BY 4.0)}
\begin{document}
\maketitle

\noindent\emph{Source and licence.} The Zero Set of an Electric Field from a Finite Number of Point Charges. Version arXiv:2106.04706v2 (CC BY 4.0), distributed by arXiv under the Creative Commons Attribution 4.0 International licence, http://creativecommons.org/licenses/by/4.0/, as recorded in the arXiv licence metadata for this exact version and shown on the abstract page https://arxiv.org/abs/2106.04706v2. Full licence notice: \texttt{licenses/arxiv-2106.04706-CC-BY-4.0.txt}.\par\smallskip

\noindent\textit{Editorial scope.} Counted unit: the article's prop countable (source0.tex lines 511--513). The article's complete original proof of it is delivered with it (source0.tex lines 515--600, one \texttt{proof} environment of 86 lines). No mathematics is written, repaired or re-derived: the statement and the proof are the article's own lines, in the article's own notation, and no retained line is substituted or re-flowed.  No further source-local context is needed: every cross-reference of every retained line resolves inside the two delivered spans.  Nothing was omitted from any retained span, so the package carries no omission record.  No retained line contains a citation, so the wrapper carries no bibliography and no bibliography entry.  No retained line contains a figure environment, an image-inclusion command or drawing code, so no image is delivered and the package declares no asset dependency.  Editorial formatting only, in addition: an AMS \texttt{amsart} preamble that restates the article's own declarations and macros used by the retained lines, namely the environments \texttt{prop} on separate counters, together with the standard packages those lines need.  The retained environments print in this document as Proposition 1 in order of appearance, whereas the article prints the same items with the numbering of its own full document.  The complete native source of the article as distributed in the arXiv e-print for this exact version is bundled unchanged as \texttt{sources/112429/source0.tex}.
\begin{prop}\label{countable}
Each point in $H_m$ is an isolated point in $H_m$ for every $m=1,2,\ldots,M-1$, unless $a_1=a_2= \cdots =a_M = 0$.
\end{prop}

\begin{proof}
Let $m$ be a fixed positive integer. Suppose $(x_0,y_0) \in H_m$, that is,
either
\begin{equation}\label{3.7} X_u(x_0,y_0) = Y_u(x_0,y_0) = 0\,, \enskip y_0 \neq 0\,, \enskip u=1,2,\ldots,m\,, \enskip X_{m+1}(x_0,y_0) \neq 0\,,
\end{equation}
or
\begin{equation}\label{3.8} X_u(x_0,y_0) = Y_u(x_0,y_0) = 0\,, \enskip y_0 \neq 0\,, \enskip u=1,2,\ldots,m\,, \enskip Y_{m+1}(x_0,y_0) \neq 0\,.
\end{equation}

If ~(\ref{3.7}) holds, then we have
\begin{equation}\label{3.9} \frac{\partial{X_m}}{\partial y}(x_0,y_0) = -(2m+1)y_0X_{m+1}(x_0,y_0) \neq 0
\end{equation}
and
\begin{equation}\label{3.10} \frac{\partial{Y_m}}{\partial x}(x_0,y_0) = -(2m+1)y_0X_{m+1}(x_0,y_0) \neq 0
\end{equation}  
Also,
\begin{equation}\label{3.11} \begin{split} \frac{\partial{Y_m}}{\partial y}(x_0,y_0) & = \, \frac{Y_m(x_0,y_0)}{y_0} - (2m+1)y_0Y_{m+1}(x_0,y_0) \cr
& = \, -(2m+1)y_0Y_{m+1}(x_0,y_0) \cr \end{split}
\end{equation}
and
\begin{equation}\label{3.12} \begin{split} \frac{\partial{X_m}}{\partial x}(x_0,y_0) & = \, (2m+1)y_0Y_{m+1}(x_0,y_0) - \frac{(2m)Y_m(x_0,y_0)}{y_0} \cr
& = \, (2m+1)y_0Y_{m+1}(x_0,y_0)\,. \end{split}
\end{equation}

By the Implicit Function Theorem and ~(\ref{3.9}), there are $\delta > 0$, $\varepsilon > 0$, and a function $y = f(x)$ differentiable on
$(x_0-\delta,x_0+\delta)$ such that the set
\begin{equation} \notag \begin{split} C_1 : & = \, \{(x,y): X_m(x,y) = 0, \enskip x \in (x_0-\delta,x_0+\delta), \enskip y \in (y_0-\varepsilon,y_0+\varepsilon)\} \cr
& = \, \{(x,f(x)): \enskip x \in (x_0-\delta,x_0+\delta)\}\cr \end{split}
\end{equation}
is the graph of a differentiable function $y = f(x)$ on $(x_0-\delta,x_0+\delta)$, and the equation for the tangent line to the curve $C_1$ at the point $(x_0,y_0)$ is
$$y-y_0 = \alpha(x-x_0)\,,$$
where
$$ \alpha = \left( \frac{\partial{X_m}}{\partial x} \Big / \frac{\partial{X_m}}{{\partial y}} \right)(x_0,y_0
= \frac{(2m+1)y_0Y_{m+1}(x_0,y_0)}{-(2m+1)y_0X_{m+1}(x_0,y_0)} = \frac{-Y_{m+1}(x_0,y_0)}{X_{m+1}(x_0,y_0)}\,.$$

Similarly, by the Implicit Function Theorem and ~(\ref{3.10}), there are $\delta > 0$, $\varepsilon > 0$, and a function $x = g(y)$
differentiable on $(y_0-\delta,y_0+\delta)$ such that the set
\begin{equation} \notag \begin{split} C_2 : & = \, \{(x,y): Y_m(x,y) = 0, \enskip y \in (y_0-\delta,y_0+\delta), \enskip x \in (x_0-\varepsilon,x_0+\varepsilon)\} \cr
& = \, \{(g(y),y): \enskip y \in (y_0-\delta,y_0+\delta)\} \cr \end{split}
\end{equation}
is the graph of a differentiable function $x = g(y)$ on $(y_0-\delta,y_0+\delta)$, and the equation for the tangent line to the curve $C_2$ at the point $(x_0,y_0)$ is
$$x-x_0 = \beta(y-y_0)\,,$$
where by ~(\ref{3.11}) and ~(\ref{3.12}) we have
$$\beta  = \left( \frac{\partial{Y_m}}{\partial y} \Big / \frac{\partial{Y_m}}{{\partial x}} \right)(x_0,y_0)
= \frac{-(2m+1)y_0Y_{m+1}(x_0,y_0)}{-(2m+1)y_0X_{m+1}(x_0,y_0)} = \frac{Y_{m+1}(x_0,y_0)}{X_{m+1}(x_0,y_0)}\,.$$

Observe that $\alpha = -\beta$. We conclude that the curves $C_1$ and $C_2$ are orthogonal to each other at $(x_0,y_0)$. As $H_m \subset C_1 \bigcap C_2$,
the point $(x_0,y_0) \in H_m$ is an isolated point of $H_m$.

If ~(\ref{3.8}) holds, then we have
\begin{equation}\label{3.13}\begin{split} \frac{\partial{Y_m}}{\partial y}(x_0,y_0) & = \, \frac{Y_m(x_0,y_0)}{y_0} - (2m+1)y_0Y_{m+1}(x_0,y_0) \cr
& = \, -(2m+1)y_0Y_{m+1}(x_0,y_0) \neq 0 \cr \end{split}
\end{equation}
and
\begin{equation}\label{3.14} \begin{split} \frac{\partial{X_m}}{\partial x}(x_0,y_0) & = \, (2m+1)y_0Y_{m+1}(x_0,y_0) - \frac{(2m)Y_m(x_0,y_0)}{y_0} \cr
& = \, (2m+1)y_0Y_{m+1}(x_0,y_0) \neq 0\,. \end{split}
\end{equation}
Also,
\begin{equation}\label{3.15} \begin{split} \frac{\partial{X_m}}{\partial y}(x_0,y_0) =  \frac{\partial{Y_m}}{\partial x}(x_0,y_0) = -(2m+1)y_0X_{m+1}(x_0,y_0) \end{split}
\end{equation}

By the Implicit Function Theorem and ~(\ref{3.13}), there are $\delta > 0$, $\varepsilon > 0$, and a function $y = f(x)$
differentiable on $(x_0-\delta,x_0+\delta)$ such that the set
\begin{equation} \notag \begin{split} C_3 : & = \,  \{(x,y): Y_m(x,y) = 0, \enskip x \in (x_0-\delta,x_0+\delta), \enskip y \in (y_0-\varepsilon,y_0+\varepsilon)\} \cr
& = \, \{(x,f(x)): \enskip x \in (x_0-\delta,x_0+\delta)\} \cr \end{split}
\end{equation}
is the graph of a differentiable function $y = f(x)$ on $(x_0-\delta,x_0+\delta)$, and the equation for the tangent line to the curve $C_3$ at the point $(x_0,y_0)$ is
$$y-y_0 = \alpha(x-x_0)\,,$$
where
$$ \alpha = \left( \frac{\partial{Y_m}}{\partial x} \Big / \frac{\partial{Y_m}}{{\partial y}} \right)(x_0,y_0)
= \frac{-(2m+1)y_0X_{m+1}(x_0,y_0)}{-(2m+1)y_0Y_{m+1}(x_0,y_0)} = \frac{X_{m+1}(x_0,y_0)}{Y_{m+1}(x_0,y_0)}\,.$$
Similarly, by the Implicit Function Theorem and ~(\ref{3.14}), there are $\delta > 0$, $\varepsilon > 0$, and a function $x = g(y)$ differentiable on
$(y_0-\delta,y_0+\delta)$ such that the set
\begin{equation} \notag \begin{split} C_4 : & = \, \{(x,y): X_m(x,y) = 0, \enskip y \in (y_0-\delta,y_0+\delta), \enskip x \in (x_0-\varepsilon,x_0+\varepsilon)\} \cr
& = \, \{(g(y),y): \enskip y \in (y_0-\delta,y_0+\delta)\} \cr \end{split}
\end{equation}
is the graph of a differentiable function $x = g(y)$ on $(y_0-\delta,y_0+\delta)$, and the equation for the tangent line to the curve $C_4$ at the point $(x_0,y_0)$ is
$$x-x_0 = \beta(y-y_0)$$
where by ~(\ref{3.15}) we have
\begin{equation} \notag \beta = \left( \frac{\partial{X_m}}{\partial y} \Big / \frac{\partial{X_m}}{{\partial x}} \right)(x_0,y_0)
= \frac{-(2m+1)y_0X_{m+1}(x_0,y_0)}{(2m+1)y_0Y_{m+1}(x_0,y_0)} = \frac{-X_{m+1}(x_0,y_0)}{Y_{m+1}(x_0,y_0)}\,.
\end{equation}

Observe that $\alpha = -\beta$. We conclude that the curves $C_3$ and $C_4$ are orthogonal to each other at $(x_0,y_0)$. As $H_m \subset C_3 \bigcap C_4$,
the point $(x_0,y_0) \in H_m$ is an isolated point of $H_m$.
\end{proof}

\end{document}
